\section{Proofs, assumptions, and the status of every statement}

The main text states each result next to the experiment it motivates. This appendix collects the statements, the assumptions under which they hold, and their proofs; where a statement is not proved, its status is said in the same place rather than left to be inferred. The summary of statuses is Table B.4.

\subsubsection{Standing assumptions}

\textbf{(A1) Exactness of the reference.} The transformed input is produced by the transformation itself --- a re-render from the transformed parameter, or an exact spectral multiplication --- never by an approximation carried out in a different colour space. Everywhere a reference route $g_\tau z$ appears, it is the image of this exact construction.

\textbf{(A2) Split disjointness.} Every fit and its evaluation use disjoint samples; where a state is said to be \textit{never fitted}, no sample from its neighbourhood was used in fitting.

\textbf{(A3) Power is reported.} Wherever $T_F$ is reported, the displacement $\lVert F(z)-F(g_\tau z)\rVert$ (the no-op error, $T_F$'s denominator) is reported or recorded with it.

\textbf{(A4) Well-definedness of $g_\tau$.} Where $g_\tau$ is claimed to be a function on activations, the fibre condition (F) below holds.

\textbf{Assumption B (linear site, invertible transformation).} On a neighbourhood of the data the site map is linear, $z=Wx$ with $W\in\mathbb R^{d\times D}$ of full row rank $d$, and $\tau$ acts through an invertible linear $g\in GL(D)$. Write $\mathcal R=\operatorname{row}(W)$ (retained), $\mathcal K=\ker(W)=\mathcal R^\perp$ (discarded), and decompose $\mathbb R^D=\mathcal R\oplus\mathcal K$ with
\[\adjustbox{max width=\columnwidth}{$\displaystyle \begin{aligned}
g&=\begin{pmatrix}A&C\\ D&B\end{pmatrix},\qquad A=P_{\mathcal R}gP_{\mathcal R},\ \ C=P_{\mathcal R}gP_{\mathcal K},\\
D&=P_{\mathcal K}gP_{\mathcal R},\ \ B=P_{\mathcal K}gP_{\mathcal K}.
\end{aligned}
$}\]

\textbf{Assumption C (isotropy).} $\mathbb E[x]=0$ and $\mathbb E[xx^\top]=\sigma^2 I_D$. This is the assumption under which the closed form of Theorem 2 is exact; the general case is Proposition B.3.

\subsubsection{The general layer (the fibre condition)}

\textbf{Observation (fibre compatibility, and the inherited composition law).} \textit{An induced action of $\tau$ on $Z_\psi$ exists if and only if} (F) \textit{holds, and when it does the induced family $\rho:M\to\operatorname{End}(Z_\psi)$ satisfies $\rho(\tau_2\tau_1)=\rho(\tau_2)\circ\rho(\tau_1)$ and $\rho(e)=I$.} This is the universal property of the quotient $Z_\psi$ written for a transformation family. \textit{Well-definedness:} if (F) holds, $\rho_\tau(\psi(s)):=\psi(\tau s)$ depends only on the class of $s$; conversely, if $\rho_\tau\psi=\psi\circ\tau$ then $\psi(s_1)=\psi(s_2)$ gives $\psi(\tau s_1)=\rho_\tau\psi(s_1)=\rho_\tau\psi(s_2)=\psi(\tau s_2)$. \textit{Inheritance:} for $z=\psi(s)$, $\rho(\tau_2)\rho(\tau_1)z=\rho(\tau_2)\psi(\tau_1s)=\psi(\tau_2\tau_1s)=\rho(\tau_2\tau_1)z$, and $\rho(e)\psi(s)=\psi(es)=\psi(s)$. $\square$

\textit{Two consequences carry the weight; the statement itself is definitional.} First, existence is a single burden rather than two --- satisfying (F) secures well-definedness and the composition law at once, so a representation either exists and composes or does not exist. Second, the composition law therefore \textbf{cannot be falsified by a network-level composition experiment}: what such an experiment measures is the defect of a fitted family, which is why \S{}5.4 grades chains against the truth-level composite and against a direct refit rather than treating them as a test of the group axiom.

\textbf{Approximate form.} If (F) holds only up to $\eta$ --- $\lVert\psi(\tau s_1)-\psi(\tau s_2)\rVert\le\eta$ whenever $\psi(s_1)=\psi(s_2)$ --- then $\rho_\tau$ is well defined on $Z_\psi$ only after quotienting by the equivalence generated by the $\eta$-ball identification, and the induced map carries a defect bounded by $\eta$. All colour and heat measurements of this paper are in the exact regime (A1); the approximate form is what the local transporter of \S{}5.3 accepts deliberately, and the composition price it pays is the one \S{}5.4 measures.

\subsubsection{The single-region linear layer (Theorems 1 and 2)}

\textbf{Lemma B.1 (the four blocks, and which one obstructs).} \textit{Write $x=x_R+x_K$ with $x_R=P_{\mathcal R}x$, $x_K=P_{\mathcal K}x$. Then}
\[\adjustbox{max width=\columnwidth}{$\displaystyle z(x)=Wx_R,\qquad z(gx)=WA\,x_R+WC\,x_K .$}\]
\textit{Proof.} $W|_{\mathcal K}=0$ gives the first identity; $gx=Ax_R+Cx_K+Dx_R+Bx_K$ and the last two terms lie in $\mathcal K$, hence are annihilated by $W$. $\square$

The blocks are not interchangeable. $A$ is the part of the action the representation can express; $D$ moves retained information into the discarded subspace and \textit{loses} information without creating a contradiction; $B$ acts inside the discarded subspace and is invisible; $C$ moves discarded information \textbf{into} the retained subspace and is, by Lemma B.1, the only term contributing to $z(gx)$ while being unavailable from $z(x)$. \textbf{$D$ costs information; $C$ costs consistency.}

\textbf{Theorem 1 (existence is kernel preservation).} \textit{Under Assumption B there exists $\rho\in\mathbb R^{d\times d}$ with $z(gx)=\rho\,z(x)$ for all $x$ if and only if $g\mathcal K\subseteq\mathcal K$, equivalently $C=0$, equivalently $\operatorname{row}(Wg)\subseteq\operatorname{row}(W)$.} \textit{Proof.} $z(gx)=Wgx$ and $\rho z(x)=\rho Wx$, so we need $\rho W=Wg$, which has a solution iff $\operatorname{row}(Wg)\subseteq\operatorname{row}(W)$. Since $\operatorname{row}(M)^\perp=\ker M$, that is $\ker W\subseteq\ker(Wg)$, i.e. $Wgx_K=0$ for every $x_K\in\mathcal K$, i.e. $g\mathcal K\subseteq\mathcal K$. By Lemma B.1 the condition is $C=0$. $\square$

\textit{Remark.} The condition constrains a \textbf{direction}, not an amount: $\lVert C\rVert$ says nothing about $\dim\mathcal K$. A rank-1 and a rank-$k$ site are equally editable when $g\mathcal K\subseteq\mathcal K$ and equally broken when $\lVert C\rVert$ is large; compression enters only through where the split is drawn.

\textbf{Theorem 2 (defect law, in the induced metric).} \textit{Under Assumptions B and C let $\rho^\star$ minimise $\mathbb E\lVert\rho z(x)-z(gx)\rVert^2$ over the population, and let $T_{RMS}$ be the RMS-criterion transfer of \S{}3.3. Then}
\[\adjustbox{max width=\columnwidth}{$\displaystyle \begin{aligned}
1-T_{RMS}&=\frac{\lVert WC\rVert_F}{\sqrt{\lVert W(A-I_{\mathcal R})\rVert_F^2+\lVert WC\rVert_F^2}}
=\frac{\kappa_W}{\sqrt{1+\kappa_W^2}},\\
\kappa_W&:=\frac{\lVert WC\rVert_F}{\lVert W(A-I_{\mathcal R})\rVert_F},
\end{aligned}
$}\]
\textit{and when $W$ is an isometry on the retained subspace this reduces to $\kappa=\lVert C\rVert_F/\lVert A-I\rVert_F$ with the optimum at $\rho^\star=A$; otherwise the optimum is $W$-conjugate.} \textit{Proof.} By Lemma B.1 the objective is $\mathbb E\lVert(\rho W-WA)x_R-WCx_K\rVert^2$. On $\mathcal R$ the operator $\rho W-WA$ acts and $x_R$ ranges over a dense subset, so the minimiser restricts to $\rho W|_{\mathcal R}=WA$; under isotropy $x_R\perp x_K$ and the normal equations decouple. The residual is $WCx_K$ and the displacement is $W(A-I)x_R+WCx_K$, so $\mathbb E\lVert WCx_K\rVert^2=\lVert WC\rVert_F^2$ and $\mathbb E\lVert W(A-I)x_R\rVert^2=\lVert W(A-I)\rVert_F^2$; dividing gives the displayed form. The mean-ratio transfer $T_F=1-\mathbb E\lVert\rho z-z'\rVert/\mathbb E\lVert z-z'\rVert$ is a different statistic and satisfies the same identity only up to the fluctuation of the norm ratio; on the sweep of Table 3 the two differ by at most $0.0011$. $\square$

\textit{Why the metric appears.} Writing the law without $W$ assumes the encoder is an isometry on the retained subspace --- true after whitening, false in general. A coordinate the encoder scales up counts more in $\lVert WC\rVert_F$, so a defect law calibrated at one site need not transfer to another. The qualitative content is metric-free: $C$ is the only block that obstructs existence (Theorem 1) and the only one that enters the defect.

\textbf{Proposition B.3 (when the closed form is tight).} \textit{For general centred $x$ with covariance $\Sigma$, the identity $\rho^\star=A$ --- hence Theorem 2 --- holds if and only if the cross term vanishes,}
\[\adjustbox{max width=\columnwidth}{$\displaystyle \begin{aligned}
&\mathbb E\big[(W(A-I)x_R)\big]^{\top}\big(WCx_K\big)=0\\
&\qquad\text{after projection on the predictable directions},
\end{aligned}
$}\]
\textit{that is, iff the $\operatorname{Cov}(x_R,x_K)$-weighted cross-covariance of the predictable part and the residual is zero. Under Assumption C it holds automatically.} \textit{Proof.} The normal equations for the minimiser contain $\mathbb E[(W(A-I)x_R)(WCx_K)^\top]$; setting $\rho=A$ satisfies them exactly when this term vanishes. $\square$

\textbf{Corollary B.4 (the score is not a property of $(W,g)$).} \textit{Under Assumption B alone, $T_F$ depends on $\Sigma$.} Verified: with $W$ and the mixing block held fixed, sweeping the anisotropy of the input distribution moves $T_F$ from $0.069$ to $0.429$ while the displacement share in the bottom-50\% variance directions moves only $0.486\to0.440$. Two consequences used in the paper: the theoretically correct instantiation of a network's discarded subspace is the \textbf{consumer-invisible} subspace rather than the low-variance directions; and cross-site $T_F$ comparisons are confounded unless the site's variance geometry is reported or matched. This is also why the law's \textit{magnitude} is stated for the controlled regime: real sites are strongly anisotropic, so Assumption C does not hold there and the closed form is not expected to be calibrated.

\textbf{Proposition B.5 (estimation-limited regime).} \textit{With $K$ samples and a $d\times d$ fit, the measured score is at least the population value minus $O(\sqrt{d^2/K})$; when $d^2/K$ is not small the measured defect over-states the irreducible one.} The two worst trained-weight cells of Table 3 (absolute error $0.23$ and $0.07$) are exactly the cells whose effective rank collapses to $4$ of $32$, which is why every reported score in this paper carries its fitted rank, its fit and held-out sample counts, and a tuned-$\lambda$ ceiling.


\begin{table*}[t]
\centering
\small
\caption*{\textbf{Table B.1.} Configuration sensitivity of the single-region law: what moves $T_F$ and by how much. Every row is a controlled sweep of one factor with the others fixed.}
\begin{tabularx}{\textwidth}{@{}>{\raggedright\arraybackslash}X >{\raggedright\arraybackslash}X >{\raggedright\arraybackslash}X >{\raggedright\arraybackslash}X@{}}
\toprule
\textbf{manipulation} & \textbf{range} & \textbf{$T_F$} & \textbf{closed-form deviation} \\
\midrule
kernel-preserving $g$ (Thm 1) & rank $8$ / $48$ & $0.9999992$ / $0.9999943$; residual $2.6\times10^{-5}$ vs displacement $34.3$ & --- \\
mixing sweep at fixed rank & $\kappa=4.5\times10^{-16}\to9.34$ & $1.0000\to0.0041$ & max $0.0091$ \\
retained-rank sweep at fixed geometry & rank $48\to4$ & $0.665\to0.341$ (factor $\approx2$) & --- \\
input anisotropy at fixed $W,\allowbreak{}g$ & $\text{aniso}=0.0\to1.0$ & $0.069\to0.429$ & --- \\
trained linear weights, 32 cells & --- & $\rho(\kappa,\allowbreak{}T_F)=-0.968$ vs $\rho(\text{eff. rank},\allowbreak{}T_F)=+0.502$ & median $0.014$ \\
\bottomrule
\end{tabularx}
\end{table*}

\textit{Reading.} Rank is not the governing variable, and the two effects are not symmetric: an arbitrarily large change in rank cannot make $\kappa$ large, whereas a small change in the geometry of the split can.

\subsubsection{b Carrier closure and optimal linear realisation}

\textbf{Proposition F (closure criterion).} \textit{For a carrier map $B$ and a carrier action $A$, there exists a fixed linear $K$ with $KB=BA$ if and only if $A\ker B\subseteq\ker B$.} \textit{Proof.} $KB=BA$ is solvable iff $BA$ vanishes on $\ker B$; a linear map is determined on a subspace only if it maps that subspace into $\mathrm{ran}\,B$. $\square$

\textbf{Proposition G (optimal linear realisation).} \textit{Let $\mathbb E[cc^\top]=I$, $z=Bc$ and $z'=BAc$. Then $K^\star=BAB^\dagger$ is a least-squares minimiser and}
\[\adjustbox{max width=\columnwidth}{$\displaystyle \min_K\mathbb E\|Kz-z'\|^2=\|BA P_{\ker B}\|_F^2,$}\]
\textit{which is zero exactly under Proposition F. For a coloured carrier, $\mathbb E[cc^\top]=\Sigma_c$, the residual is $\|BA\Sigma_c^{1/2}P_{\ker(B\Sigma_c^{1/2})}\|_F^2$ and the minimiser is $K^\star=BA\Sigma_c B^\top(B\Sigma_cB^\top)^\dagger$, with $\Sigma_c^{1/2}$ the symmetric square root: correlation lets a retained coordinate predict a discarded one, so the projector is taken after whitening, and the two forms coincide only for $\Sigma_c\propto I$.} \textit{Proof.} $K^\star z-z'=BA(B^\dagger B-I)c=-BAP_{\ker B}c$ because $B^\dagger B$ is the orthogonal projector onto the row space of $B$; taking the expectation with $\mathbb E[cc^\top]=I$ gives the squared Frobenius norm, and the coloured case follows by whitening $c\mapsto\Sigma_c^{-1/2}c$. $\square$

\textit{Corollary (coordinate selection).} If $B=P$ selects coordinates and the carrier is orthogonal along the orbit, $K^\star=PAP^\top$: the optimal fixed linear map is the carrier action projected onto the retained coordinates. Verified over forty random selections of a six-dimensional carrier, where the gap to $PAP^\top$ was $9\times10^{-16}$ and the closure verdict coincided with the sign of the residual in every case.

\textit{Corollary (conjugate pairs).} For a full continuous rotation action written in one two-dimensional block per frequency, with selection along those coordinates, closure holds iff each block is retained or dropped as a whole. The restriction is part of the statement: for a single fixed angle, or when frequencies repeat and their blocks may mix, the relevant object is the invariant subspace rather than the pair. The instantiation on measured features, including the dropped-coordinate variant, is Table D.21.

\subsubsection{The piecewise-linear layer (Propositions C, D, E)}

A rectifier network partitions its input space into \textbf{activation regions} $R_G$ indexed by the gate pattern $G(x)=\mathbf{1}[a(x)>0]$, and on each region the site map is affine, $\phi(x)=A_Gx+b_G$. The gate is recoverable from the site tensor itself, since $z=\mathrm{ReLU}(a)$ gives $G=\mathbf{1}[z>0]$. The \textbf{crossing set} is $\mathcal C_\tau=\{x:G(\tau x)\neq G(x)\}$.

\textbf{Proposition C (region-transition criterion).} \textit{Let the site be affine on each visited region, $\phi=A_rx+b_r$, let $T$ be the input transformation, and let $(r,s)$ range over the realised transitions, each covering a set with non-empty interior or spanning its region. Then a fixed linear $\rho$ with $\rho\phi=\phi\circ T$ on the visited cells exists if and only if the joint system $\rho A_r=A_sT,\ \rho b_r=b_s$ is consistent over all realised $(r,s)$.} \textit{Proof.} On a cell both $\rho\phi(x)$ and $\phi(Tx)$ are affine in $x$; an identity between affine maps on a set with non-empty interior extends to the maps, so the system is necessary, and a solution clearly suffices. $\square$

\textbf{Corollary C.1 (no crossing).} \textit{If $T$ leaves the gate pattern of each visited region unchanged, the realised transitions are the diagonal cells and the system reduces to $\rho A_r=A_rT,\ \rho b_r=b_r$.} The special case is the one in which each source region maps into itself; consistency with crossings is independent of it, as the two counterexamples of \S{}3.5 show.

For a single visited region the system is $\rho A=A\tau,\ \rho b=b$, which is Theorem 1's kernel condition rewritten --- the linear layer is the single-region case of this one.

\textit{Consequence.} \textbf{Even with no crossing at all, the multi-region structure alone can obstruct a global linear action}: the visited regions must be simultaneously conjugate to the transformation. This is why enlarging the operator family buys little --- measured, a ridge-plus-$3\times3$-residual family differs from the global linear family by at most $0.087$ at every ResNet-50 site and $0.11$ at every ConvNeXt site, with a median gap of $0.005$ over the 32-site grid; the two separate only where both fail (DINOv2 heat $block8$: $-0.40$ against $-0.99$).

\textbf{Proposition D (the decomposition, and where rule diversity comes from).} \textit{With $z=\mathrm{ReLU}(a)$, $a'=a+\Delta a$, $G=\mathbf{1}[a>0]$ and $G'=\mathbf{1}[a'>0]$,}
\[\adjustbox{max width=\columnwidth}{$\displaystyle z'-z=G\odot(a'-a)+(G'-G)\odot a',$}\]
\textit{and on a unit whose gate changes the second term is $+a'$ (off$\to$on) or $-a'$ (on$\to$off).} \textit{Proof.} $z'=G'\odot a'=G\odot a'+(G'-G)\odot a'$ and $z=G\odot a$. $\square$

Two consequences, and the second governs how the measurements are read. First, the within-regime term is linear in $\Delta a$; the gate-change term has magnitude $|a'|$ on the changed units, so it is not a constant that survives a shrinking perturbation. Second, \textbf{a gate change is not itself an obstruction}: it is the origin of a change in the local affine rule, and whether a single fixed linear operator can satisfy all the rules the transformation visits is decided by the joint system of \S{}3.5, which may be consistent with every gate changed, or inconsistent with no gate changed in the source region. Both directions have explicit two-dimensional counterexamples in \S{}3.5. Measured (script 247; ResNet-50, three depths $\times$ two algebras): $\rho(\text{flip count},T_F)=+0.49$ (n.s.) against $\rho(\text{flip depth},T_F)=\rho(\text{gate-change share},T_F)=-0.83$.

\textbf{Proposition E (achievable score of the fixed-linear class).} \textit{Let $F$ be the consumer, $\lVert\cdot\rVert_F$ the seminorm its response induces, and $d=z'-z$. For a linear consumer the identity}
\[\adjustbox{max width=\columnwidth}{$\displaystyle T_F^{\max}=1-\frac{\lVert\mathrm{Proj}_F d^\perp\rVert_F}{\lVert\mathrm{Proj}_F d\rVert_F}$}\]
\textit{is exact, with $d^\perp$ the residual of the $L^2$ projection of $d$ on the consumer-visible linear span of $z$: the class of fixed linear maps can reproduce any component of $d$ that is linearly predictable from $z$, and nothing else.} \textit{Proof.} Write $d=d^{\parallel}+d^\perp$ with $d^{\parallel}$ in the visible span of $z$; the map $K=d^{\parallel}$ is attainable and leaves the residual $d^\perp$, which no member of the class can reach, so the optimum is $d^{\parallel}$ and the displayed ratio is the score attained. $\square$

\textit{Status.} Proved for linear consumers. No claim is made for families outside the fixed-linear class; the proxy is evaluated at six (site, algebra) points in \S{}6.3.

\textit{The gate-change share is a proxy.} With $\Delta_{\mathrm{flip}}:=(G'-G)\odot a'$ the gate-change term of Proposition D, its consumer-visible share approximates the slack when that term is the part of $d$ not linearly predictable from $z$. It is an identity in neither direction: the all-gates-change counterexample of \S{}3.5 has zero slack, and a consumer blind to the change has zero visible share whatever the slack. The six-point check is the paper's evidence that the proxy is informative at the measured sites; the two-source decomposition of \S{}3.3 is the statement about the failure itself.

\subsubsection{Realisation and composition (Theorems 3, 4, 5)}

\paragraph{B.5.1 Coordinate-monotone families on a closed orbit}

\textbf{Definition B.6.} A map $G:\mathbb R^d\to\mathbb R^d$ is \textit{coordinate-wise monotone} if $G(z)_i=f_i(z_i)$ with each $f_i:\mathbb R\to\mathbb R$ non-decreasing. A family $\{G_t\}_{t\in I}$ is a \textit{monotone parameter family} if every $G_t$ is coordinate-wise monotone and, for every $i$, the map $(t,u)\mapsto f_{i,t}(u)$ is non-decreasing in both arguments. Such a family \textit{realises} the orbit $O=\{z(t)\}$ from the anchor $z(0)=z^0$ if $G_t(z^0)=z(t)$ for all $t\in I$.

\textbf{Lemma B.7 (obstruction).} \textit{If a monotone parameter family realises $O$ from $z^0$, then every coordinate $t\mapsto z(t)_i$ is non-decreasing. Consequently, if $O$ is closed ($z(1)=z(0)$) and has a coordinate that decreases somewhere, no monotone parameter family realises it exactly.} \textit{Proof.} $z(t)_i=f_{i,t}(z^0_i)$ and $t\mapsto f_{i,t}(u)$ is non-decreasing for fixed $u=z^0_i$. $\square$

\textbf{Lemma B.8 (quantitative one-coordinate cost).} \textit{Let $g:[0,1]\to\mathbb R$ with $g\ge\alpha$ on an interval $I$ and $g\le\alpha-A$ on a later interval $J$, $A>0$, and let $|I|\le|J|$. Then for every non-decreasing $c$,}
\[\adjustbox{max width=\columnwidth}{$\displaystyle \int_0^1(g-c)^2\;\ge\;\frac{A^2}{4}\,|I| .$}\]
\textit{Proof.} Suppose $\int_I(g-c)^2<\frac{A^2}{4}|I|$ and $\int_J(g-c)^2<\frac{A^2}{4}|J|$. On $I$ and $J$ respectively the sets where $c<\alpha-A/2$, resp. $c>\alpha-A/2$, then have measure less than $|I|/2$, resp. $|J|/2$. So there is $t_0\in I$ with $c(t_0)\ge\alpha-A/2$ or $u_0\in J$ with $c(u_0)\le\alpha-A/2$. In the first case monotonicity gives $c\ge\alpha-A/2$ on all of $J$ (which lies to the right of $I$), so $(g-c)^2\ge A^2/4$ on $J$ and the second integral is at least $\frac{A^2}{4}|J|\ge\frac{A^2}{4}|I|$, a contradiction; in the second case $(g-c)^2\ge A^2/4$ on $I$ by the same argument with the roles exchanged. $\square$

\textbf{Theorem 3 (a closed orbit has no exact coordinate-monotone realisation).} \textit{A family $\{G_t\}$ whose members are coordinate-wise monotone in a fixed basis, with $(t,u)\mapsto f_{i,t}(u)$ non-decreasing in both arguments, realises an orbit only if every coordinate of that orbit is monotone in the parameter (Lemma B.7); consequently a closed non-degenerate orbit admits no exact realisation by such a family. Quantitatively, every coordinate that decreases by $A$ over intervals of length at least $\ell$ contributes $L^2$ error at least $A^2\ell/4$ (Lemma B.8).} The empirical content is the comparison in Table 7: the class is rejected both by the theorem and by its measured best fit on the reference orbit. \textit{Scope.} The statement is about monotone parameter families, which is what the lemmas prove; an aggregate bound in terms of the third principal value of the orbit covariance does not follow from them and is not claimed here.

\paragraph{B.5.2 Accumulation law}

\textbf{Theorem 4 (accumulation law).} \textit{Let the truth action $\tau$ be $L$-Lipschitz in the consumer's seminorm on the consumed range, and let a realised family satisfy the uniform single-step bound $\lVert\rho(\tau)z-g_\tau z\rVert_F\le\varepsilon$ on the states visited by the chain. Then}
\[\adjustbox{max width=\columnwidth}{$\displaystyle \begin{aligned}
E_n&\le\varepsilon+L\,E_{n-1},\qquad\text{so}\\
E_n&\le\varepsilon\,\frac{L^n-1}{L-1}\ (L\neq1),\qquad E_n\le n\varepsilon\ (L=1).
\end{aligned}
$}\]
\textit{Proof.} Insert and subtract the truth's intermediate state:
\[\adjustbox{max width=\columnwidth}{$\displaystyle \begin{aligned}
\rho(\tau)^nz-g_{\tau^n}z
&=\big[\rho(\tau)\big(\rho(\tau)^{n-1}z\big)-g_\tau\big(\rho(\tau)^{n-1}z\big)\big]\\
&\quad+\big[g_\tau\big(\rho(\tau)^{n-1}z\big)-g_\tau\big(g_{\tau^{n-1}}z\big)\big].
\end{aligned}
$}\]
The first bracket is at most $\varepsilon$ by the uniform single-step bound applied at the state $\rho(\tau)^{n-1}z$; the second is at most $L\lVert\rho(\tau)^{n-1}z-g_{\tau^{n-1}}z\rVert_F=L\,E_{n-1}$ by Lipschitzness. Iterating gives the closed forms. $\square$

\textit{Uniformity.} Both hypotheses must hold uniformly over the states the chain visits. If the single-step error or the Lipschitz constant varies with the state, the recursion holds pointwise with the local constants and the closed forms become the corresponding products. Estimating $\varepsilon$ and $L$ independently of the measured chain error, rather than fitting the recursion to it, is not done here; which of the three regimes a site occupies is measured in the consumer's seminorm (Table D.22), not inferred from the generator.

\paragraph{B.5.3 Carrier sufficiency}

\textbf{Setting.} At the site, let the carrier be $c$, the synthesis $B$, the target $z'=BAc+w'$ with $\lVert w'\rVert_{L^2}\le\delta'$, the read-in $E_\theta$ (any measurable map, in particular the learned MLP), and let the interface apply $A_{\mathrm{int}}$ and decode by $B$. Band-limitation and plane sharing are the two premises the construction rests on and both are measurable, but the theorem below does not need them as hypotheses: it is an identity plus a triangle inequality for whatever $B$, $A$ and $E_\theta$ are.

\textbf{Theorem 5 (carrier sufficiency, three-term bound).} \textit{Write $e(z)=E_\theta(z)-c$ for the read-in error. Then}
\[\adjustbox{max width=\columnwidth}{$\displaystyle BA_{\mathrm{int}}E_\theta(z)-z'=BA_{\mathrm{int}}\,e(z)+B(A_{\mathrm{int}}-A)c-w'$}\]
\textit{holds exactly, and hence}
\[\adjustbox{max width=\columnwidth}{$\displaystyle \begin{aligned}
\big\|BA_{\mathrm{int}}E_\theta(z)-z'\big\|_{L^2}\le{}&\|BA_{\mathrm{int}}\|\,\|e\|_{L^2}\\
&+\|B(A_{\mathrm{int}}-A)\|\,\|c\|_{L^2}+\delta'.
\end{aligned}
$}\]
\textit{Proof.} Add and subtract $BA_{\mathrm{int}}c$ and $BAc$; the first term is the read-in error, the second the action mismatch, the third the target's tail. The triangle inequality gives the bound, with no cross-term cancellation assumed. $\square$

\textit{Three consequences.} The bound applies to the interface as built: $E_\theta$ need not be linear, and a linear read-in is the special case in which $\lVert e\rVert$ is written as an operator norm. Its three terms are read-in error, action mismatch and target tail; only the first is removed by training, and the second vanishes when the interface uses the true carrier action. The mismatch term is first order in $\lVert A_{\mathrm{int}}-A\rVert$, so the squared error is second order in the plane misalignment, while the measured quantity is the error norm, whose log-log slope in the tilt angle is $0.99$. If the carrier tail is normalised as an energy fraction, $\mathbb E\lVert w\rVert^2\le\varepsilon_{\mathrm{tail}}^2\mathbb E\lVert z\rVert^2$, the measured $12$--$16\%$ energy fraction corresponds to $\varepsilon_{\mathrm{tail}}\approx0.35$--$0.40$, not to $0.12$--$0.16$. The closure obstruction of B.3b enters separately and constrains a \textit{fixed linear map on the measured features}, a class the learned, possibly nonlinear read-in escapes; what it cannot escape is the target's tail and any mismatch between the action it applies and the action the carrier has.

\subsubsection{Status of every statement}


\begin{table*}[t]
\centering
\small
\caption*{\textbf{Table B.4.} Statement status. \textit{Definition} = fixes vocabulary; \textit{Proved} = proof above; \textit{Proved + verified} = proof plus a controlled numerical check; \textit{Empirical} = a measured bound or law, not a theorem; \textit{Hypothesis} = stated as such in the text.}
\begin{tabularx}{\textwidth}{@{}>{\raggedright\arraybackslash}X >{\raggedright\arraybackslash}X >{\raggedright\arraybackslash}X >{\raggedright\arraybackslash}X@{}}
\toprule
\textbf{statement} & \textbf{status} & \textbf{scope} & \textbf{where} \\
\midrule
Definition 1 (representation of the algebra) & Definition & general encoders and algebras & \S{}3.1 \\
Definition 2 (representation vs behavioural level) & Definition & general & \S{}3.1 \\
Observation (fibre compatibility and inherited composition) & Proved (one line) & general & B.2 \\
Theorem 1 (kernel preservation) & Proved + verified & linear sites, Assumption B & B.3 \\
Theorem 2 (defect law) & Proved + verified & Assumptions B, C & B.3 \\
Proposition B.3 (tightness under general covariance) & Proved & Assumption B & B.3 \\
Corollary B.4 (score depends on $\Sigma$) & Proved + verified & Assumption B & B.3 \\
Proposition B.5 (estimation-limited regime) & Proved & finite samples & B.3 \\
Proposition C (region-transition criterion) & Proved & piecewise-linear sites, realised transitions & B.4 \\
Proposition D (decomposition and rule diversity) & Proved & rectifier sites & B.4 \\
Proposition E (exact score of the fixed-linear class) & Proved for linear consumers; the gate-change share is an empirical proxy for its slack & fixed-linear class & B.4 \\
Theorem 3 (no exact coordinate-monotone realisation of a closed orbit) & Proved & monotone parameter families in a fixed basis & B.5.1 \\
Theorem 4 (accumulation law) & Proved; the regime at a site is measured in the consumer's seminorm & any realised family, uniform single-step bound & B.5.2 \\
Proposition F (closure criterion), Proposition G (optimal linear realisation) & Proved + verified & carrier and action; any site & B.3b \\
Theorem 5 (carrier sufficiency, three-term bound) & Proved (identity) & any carrier, action and read-in & B.5.3 \\
Harmonic concentration, shared scaffold, inheritance & Measured & colour instance & \S{}4 \\
Depth erosion follows the readout--transformation pairing & Measured & four backbones, real images & \S{}6, Table 11 \\
Learning moves $T_F$ only through $\Delta\kappa$ & Measured (linear control) & controlled linear setting & \S{}6.4, Table 12 \\
Interface reads hue and edits regions & Measured & colour instance & \S{}7, Table 15 \\
\bottomrule
\end{tabularx}
\end{table*}

